% ==============================================================================
% sections/conclusion.tex
% Created: 2026-05-08
% ==============================================================================

Classification error in survey responses to employment questions is pervasive
and, even when small, systematically inflates measured employment transition
rates. We have developed a structural EM-algorithm estimator that jointly
identifies true transition rates and misclassification probabilities from
three consecutive waves of panel data, requiring no reinterview data or
administrative validation. The estimator exploits the fact that transient
misclassification causes the observed employment sequence to violate the
first-order Markov property in a specific, quantifiable way, and converts
this insight into a fast, numerically reliable sequence of closed-form
parameter updates.

Applied to the South African Quarterly Labour Force Survey, the estimator
yields a misclassification probability of approximately 3\%, which reduces
estimated quarterly entry and exit rates from roughly 8\% to roughly 3\%---a
correction factor of nearly three. This finding is robust across all
specifications we consider: controlling for observable worker characteristics
(age, education, gender, race, contract type) reduces the estimated
misclassification probability moderately but does not eliminate it; allowing
unobserved two-type heterogeneity in transition rates leaves the misclassification
estimate essentially unchanged while resolving an implausible ``coin-flip''
high-churn type; extending the Markov order to AR(2) with four waves reveals
substantial duration dependence but retains a positive misclassification
probability; and the tenure contamination model provides an independent
structural corroboration using job tenure data. Survey respondents with
objectively inconsistent age or education progressions exhibit employment
dynamics consistent with misclassification rates more than fifteen percentage
points higher than reliable respondents, providing direct mechanism evidence
that the estimated misclassification reflects genuine survey error.

The policy implications are substantial. If true quarterly entry and exit rates
are approximately 3\%, the average employment spell in South Africa is on the
order of several years---considerably longer than the raw transition matrix
would suggest. This characterisation of South African employment as highly
persistent implies that unemployment reflects primarily a failure of job creation
rather than a failure of job retention. The policy prescription is correspondingly
different: interventions that lower barriers to hiring, promote labour demand,
and reduce the cost of formal employment are more likely to be effective than
policies focused on worker protections, employment insurance, or the facilitation
of job search. This does not imply that active labour market policies are
irrelevant---workers who are persistently non-employed may still benefit from
training and placement programmes---but it does imply that the design of such
programmes should be predicated on a realistic understanding of the degree of
labour market dynamism.

Beyond South Africa, our estimator is applicable wherever three-wave panel data
on a binary state are available and reinterview data are not. The measurement
error problem we address is particularly acute in developing country surveys,
where reinterviews are uncommon and measurement quality may be lower, but it
arises in any survey context where respondents may conflate employment states
or where the survey instrument introduces transient coding errors. Potential
extensions include applications to poverty transitions \citep{Bigsten2008},
health insurance coverage \citep{van2017critique}, and educational enrolment
\citep{fredriksson1997economic}, all of which involve binary state variables
measured with potential transient error in panel surveys. The misclassification
problem is not unique to employment, and the EM approach extends naturally to
any setting where the first-order Markov assumption can be motivated for the
underlying true process.

Several limitations of our approach merit acknowledgment. First, the baseline
misclassification structure assumes symmetric and time-invariant error
probabilities. The inconsistency-augmented specification relaxes the
homogeneity assumption, but the asymmetric baseline results suggest that the
symmetry restriction is innocuous in our data. Second, we have treated the
first-order Markov assumption on the latent true process as the baseline
identification assumption, relaxing it only in the AR(2) extension. Higher-order
dynamics, time-varying transition rates, or more complex dependence structures
could in principle be accommodated within the EM framework at the cost of
additional data requirements. Third, our approach identifies misclassification
from the time-series variation in employment status across waves; it cannot
separately identify misclassification that is perfectly correlated across waves
(a worker who is always misclassified would appear to have a different true
state, not a noisy version of their actual state). Future work could extend
the approach to longer panels, time-varying misclassification probabilities,
or multivariate employment state variables (employed, unemployed, inactive),
all of which are within the scope of the EM framework developed here.
